%
% $Header: /home/cvs/root/haskell-report/report/standard-prelude.verb,v 1.4 2002/12/10 11:51:11 simonpj Exp $
%
%**<title>The Haskell 98 Report: Standard Prelude</title>
%**~header
%*anchor on
\chapter{Standard Prelude}
\label{stdprelude}

In this chapter the entire \Haskell{} Prelude is given.  It constitutes
a {\em specification} for the Prelude.  Many of the definitions are
written with clarity rather than efficiency in mind, 
and it is not required that the specification be implemented as shown here.

The default method definitions, given with \mbox{\tt class} declarations, constitute
a specification {\em only} of the default method.  They do not constitute a
specification of the meaning of the method in all instances.  To take
one particular example, the default method for \mbox{\tt enumFrom} in class \mbox{\tt Enum}
will not work properly for types whose range exceeds that of \mbox{\tt Int} (because
\mbox{\tt fromEnum} cannot map all values in the type to distinct \mbox{\tt Int} values).

The Prelude shown here is organized into a root module, \mbox{\tt Prelude},
and three sub-modules, \mbox{\tt PreludeList}, \mbox{\tt PreludeText}, and \mbox{\tt PreludeIO}.
This structure is purely presentational.
An implementation is not required to use 
this organisation for the Prelude,
nor are these three modules available for import separately.
Only the exports of module \mbox{\tt Prelude} are significant.

Some of these modules import Library modules, such as \mbox{\tt Data.Char}, \mbox{\tt Control.Monad}, \mbox{\tt System.IO},
and \mbox{\tt Numeric}.  These modules are described fully in Part~\ref{libraries}.
These imports are not, of course, part of the specification
of the \mbox{\tt Prelude}.  That is, an implementation is free to import more, or less,
of the Library modules, as it pleases.

Primitives that
are not definable in \Haskell{}, indicated by names starting with ``\mbox{\tt prim}'', are
defined in a system dependent manner in module \mbox{\tt PreludeBuiltin} and
are not shown here.  Instance declarations that simply bind primitives to
class methods are omitted.  Some of the more verbose instances with
obvious functionality have been left out for the sake of brevity.

Declarations for special types such as \mbox{\tt Integer}, or \mbox{\tt ()} are
included in the Prelude for completeness even though the declaration
may be incomplete or syntactically invalid.  An ellipsis ``\mbox{\tt ...}'' is often
used in places where the remainder of a definition cannot be given in Haskell.

To reduce the occurrence of unexpected ambiguity errors, and to
improve efficiency, a number of commonly-used functions over lists use
the \mbox{\tt Int} type rather than using a more general numeric type, such as
\mbox{\tt Integral\ a} or \mbox{\tt Num\ a}.  These functions are: \mbox{\tt take}, \mbox{\tt drop},
\mbox{\tt !!}, \mbox{\tt length}, \mbox{\tt splitAt}, and \mbox{\tt replicate}.  The more general
versions are given in the \mbox{\tt Data.List} library, with the prefix
``\mbox{\tt generic}''; for example \mbox{\tt genericLength}.

\clearpage
% The index entries for :, [], (), and tuples are here
% it just so HAPPENS that they'll end up referring to the right page
% HHAACCKK!!   <---- Partain, you scoundrel! -- KH
% Well, they aren't there anymore so I nuked them.  So there!  jcp
\inputHS{Prelude}\clearpage

%\section{Prelude {\tt PreludeBuiltin}}
%\label{preludebuiltin}
%\input{PreludeBuiltin}\clearpage


\section{Prelude \texorpdfstring{{\tt PreludeList}}{PreludeList}}
\label{preludelist}
\inputHS{PreludeList}\clearpage
\section{Prelude \texorpdfstring{{\tt PreludeText}}{PreludeText}}
\label{preludetext}
\inputHS{PreludeText}\cleardoublepage

\section{Prelude \texorpdfstring{{\tt PreludeIO}}{PreludeIO}}
\label{preludeio}
\inputHS{PreludeIO}
%**~footer


